\section{Non-Technical Summary}

When hospitals use computer programs to help predict patient outcomes  such as who might respond well to a treatment, or who is at highest risk of deteriorating  those programs are usually tested on data from the same hospital or study that created them. The problem is that patients at a different hospital, in a different part of the country, or from a different background may not look like the patients the program learned from. When that happens, the program's predictions can become unreliable, sometimes performing no better than a random guess.
Research published in Science in 2024 showed this happening clearly: models that appeared highly accurate within their original clinical trials dropped to coin-flip performance when tested on independent trials. This is not a rare or unusual failure  similar patterns have been found across genomics, immunology, and biomarker research. The tools to explain why this happens already exist. What does not exist is a practical way to check, before a program is used on real patients, whether it is likely to work in that specific setting.
This project addresses that gap. The idea is straightforward: before deploying a prediction program on a new group of patients, measure how different those patients are from the ones the program was trained on. If the difference is small, the program is probably safe to use. If the difference is large, it probably is not. The framework also checks which pieces of information the program relies on  some are reliable across different patient groups, while others only hold for the specific group the program was built on.
The result is a simple scoring system and decision flowchart that tells clinicians and data scientists whether to go ahead, make adjustments, or hold back. Early testing on real ICU data from 186 US hospitals shows the approach can flag high-risk deployment scenarios. It does not solve every problem in medical AI, but it addresses the most dangerous one: using a tool that looks good on paper without knowing whether it will actually help the patients in front of you.
